\documentclass[11pt]{article}
% Shared preamble for the round 1 response letters.
% RED-D-26-00078, submitted as "Housing First or Treatment First?",
% resubmitted as "Housing First or Treatment First? Ask the Kids, Not the Parents"
%
% House style, decided once here:
%   * responses are in RED, referee text is in the black italic block above them
%   * NO HYPHENS anywhere in the prose
%   * every response ends with a pointer to where in the new paper it lives

\usepackage[margin=1.05in]{geometry}
\usepackage{amsmath,amssymb}
\usepackage{xcolor}
\usepackage[official]{eurosym}
\usepackage{booktabs}
\usepackage{longtable}
\usepackage{array}
\usepackage[hidelinks]{hyperref}

% enumitem is not in this TeX installation, so first level list spacing is set
% directly instead
\makeatletter
\renewcommand{\@listI}{%
  \leftmargin=1.7em \labelwidth=1.2em
  \topsep=0.4em \parsep=0pt \itemsep=0.3em \partopsep=0pt}
\makeatother

\definecolor{respred}{RGB}{190,25,25}

\setlength{\parskip}{0.35em}
\setlength{\parindent}{0pt}

\newcolumntype{L}[1]{>{\raggedright\arraybackslash}p{#1}}
\newcommand{\pname}[1]{\texttt{#1}}

% ---- the two building blocks -------------------------------------------------
% \cmt{label}{short title} opens a comment, then the verbatim or paraphrased
% referee text goes in the referee environment, then the response.

\newcommand{\cmt}[2]{%
  \par\medskip
  \noindent\textbf{#1\quad #2}\par}

\newenvironment{referee}
  {\begin{list}{}{\leftmargin=1.6em \rightmargin=1.2em \topsep=0.4em}\item[]\small\itshape}
  {\end{list}}

\newenvironment{response}
  {\par\begingroup\color{respred}\medskip}
  {\par\endgroup\medskip}

% a flag for statements in the reports that are factually wrong, so they are
% easy to find again and easy to handle carefully
\newcommand{\correction}{\textbf{[Point of correction.]}\ }
\newcommand{\concede}{\textbf{[Conceded.]}\ }
\newcommand{\newresult}{\textbf{[New result.]}\ }
\newcommand{\open}{\textbf{[Still open.]}\ }

% numbers that are not in yet, so they are impossible to miss on a read through
\newcommand{\tbd}[1]{\textbf{[TO FILL: #1]}}

% Marks what has been done since the first draft of these letters. The letters
% were written when most of this was a plan; this flag marks the places where it
% is now built and running.
\newcommand{\thisround}{\textbf{[Built since the first draft of this letter.]}\ }

\newcommand{\lettertitle}[2]{%
  \begin{center}
  {\Large\bfseries Housing First or Treatment First?}\\[0.3em]
  {\normalsize Manuscript RED-D-26-00078}\\[0.9em]
  {\large\bfseries #1}\\[0.3em]
  {\normalsize #2}
  \end{center}
  \vspace{0.6em}}


\begin{document}

\lettertitle{Response to Referee 2}{Conceptual and detail concerns, responses in red}

Thank you for the report. Of the two welfare margins in your second paragraph, the
second, what becomes of the children of drug users, is now the centre of the paper,
although it enters the model differently from how you framed it (C3). I answer the
first at C2.


{\color{respred}
\textbf{One note before the comments.} The paper has a new title: it was submitted as
\emph{Housing First or Treatment First?} and is now \emph{Housing First or Treatment First? Ask the Kids, Not the Parents}. The reason is the substance of the revision rather than a relabelling, and it
will be clearest at C3, which is about the children the title refers to.

\textbf{What changed, in brief.} The paper now asks whether housing assistance should
require sobriety, within the rationed voucher programme that exists, and answers it with
a calibrated general equilibrium life cycle model in which a parent's drug use harms its
child and the parent values its child only in part. Moral hazard alone does not justify
the condition. Without the externalities on children, a newborn is worse off under
Treatment First, by 1.09 percent of lifetime consumption, even when the savings from the
drug test buy vouchers for other households. With them, in the long run and with the drug
price and the chance of drawing a voucher both adjusting, Treatment First cuts drug use by
21 percent and homelessness by 11 percent, and a newborn gains 2.57 percent of lifetime
consumption, three quarters of it an externality its parent does not count. The parents
gain too, by 0.69 percent of their own (Sections 6.2 and 6.3, Tables 5 and 6). Along the
main transition, neither homelessness nor the share of children in a shelter rises
(Section 6.4).

Several changes bear directly on your report. The children of users, whom you raised in
C3, are at the centre of the paper: a parent counts only part of what its drug use costs
its child, and the rest is an externality. The effect of housing on recovery from a
substance use disorder (SUD) is calibrated to Milby et al.\ (2005), and the model comes
close to the same trial's comparison of contingent with unconditional housing without being
targeted; the response of drug use to the condition itself is set to match the US
military's random drug testing (C1). Preferences are the same for every household, whatever
its disorder, its mental health or its dwelling, and a new two period stylized model
(Section 3) shows that it is the uninsurable risk of dying, not the utility function, that
puts drug use at the bottom of the income distribution (C4, C5). Homeownership, house
prices, endogenous rents and labour supply are gone, and the housing grid is a shelter and
three rental sizes (D1). Four things you asked for are not in the paper, or are only partly
in it, and I say so where they come up: the disamenity of homelessness is a robustness
check rather than part of the benchmark, and Corbae, Glover and Nattinger (2024) is not
cited (C2); the paper does not plot the drug use policy function by income and wealth (C4);
it does not solve for the optimal voucher programme or try broader eligibility or a more
generous voucher (C6); and it does not test a uniformly spaced rental ladder (D1).
}

%======================================================================
\section*{Conceptual concerns}

\cmt{C1}{The welfare results rest on the elasticity of drug use to housing access}
\begin{referee}
``\dots all of the welfare results rest on whether the model has an accurate
elasticity of drug use to housing access. The only empirical drug use elasticity
mentioned is a price elasticity of drug use participation, but it isn't clear to
me that this provides enough or the correct discipline for these policy
experiments. I would like to see that the model matches some evidence on the
response of drug use to housing access (see Milby, et al (2005) for cocaine users,
but there are many papers on this topic.)''
\end{referee}
\begin{response}
You are right that a price elasticity cannot by itself discipline a housing policy. The
version you read targeted the participation elasticity; the revised calibration targets
neither price elasticity, and both are now untargeted checks. The Associate Editor also
named the elasticity of drug use to housing access as one the model should match. The
response of drug use to housing, and to a condition on housing, is now tied to evidence
at three points, two of them from Milby et al.\ (2005), the trial you suggested.

\textbf{Housing and recovery.} A household with a SUD is less likely to recover in a
shelter than in a home, by the multiplier $\mu_H$ on the recovery hazard (Section 4.6,
equation (18)). $\mu_H$ is calibrated to Milby et al.: among homeless adults with a
cocaine disorder in treatment, those who were housed had 1.47 times the abstinence of those
who were not, and the model gives 1.45 (Section 5.1, Block 4, and Table 3). Housing changes
the chance of recovering, not the preference parameters for the drug, which also bears on
your C5. Housing also makes life worth more, and a household with more to lose from dying
is deterred more by the mortality risk of a SUD (Section 5.1, Block 2).

\textbf{The response to the condition itself.} The experiments also depend on how much drug
use responds when the voucher is made conditional on not using, which is a response to
housing access rather than to the price. The detection rate of Treatment First's screen is
set so that the screen deters use as much as a drug test with a severe penalty does in the
closest evidence. After the US military began testing its personnel at random and
dismissing those who failed, service members who stayed were 61 percent less likely to have
used marijuana in the past year (Mehay and Pacula, 1999). In the model, with a detection
rate of 20 percent, 62 percent of the users who face the screen stop when it is introduced
(Section 6, ``The screen''). The two are compared in proportion, since the military figure
counts any use of marijuana (footnote 34). The 61 percent is the raw comparison among those
who stayed, and the decline is 11 to 45 percent once those who left are counted (footnote
33), so I also run the screen at half the rate. Then 41 percent of users stop, inside that
range. In the long run Treatment First then cuts drug use by 14 percent and homelessness by
6 percent, and a newborn still gains 1.65 percent of lifetime consumption, against 2.57 at
the calibrated rate (footnote 35).

\textbf{Two untargeted checks.} Milby et al.\ also randomised housing that was contingent on
abstinence against housing that was not. The contingent arm reached 8.32 consecutive weeks
of abstinence against 6.60, a ratio of 1.26. The model's counterpart, the recovery hazard of
voucher holders with a SUD under Treatment First against that of housed households with a
SUD and no voucher, is 1.40 in the first year of the condition, the horizon of the trial
(Section 5.2, Block 8). The same trial therefore calibrates the effect of housing and checks
the effect of the condition. Costantini (2025) adds a check on mortality rather than on
drug use: enrolling in the Housing First programme of the Department of Veterans Affairs was
estimated to lower the three year mortality of homeless veterans with mental illness by 4.6
percentage points. In the model, with the drug price and the chance of drawing a voucher
held, households with serious mental illness (SMI) whom Treatment First's revocation puts on
the street have a three year mortality about 3.6 percentage points higher than had they kept
their voucher, the same order of magnitude (Section 5.2, Block 7). The paper compares the
model with his Housing First estimate only, not with his finding on the Department's
Treatment First programmes.

\textbf{The wider literature.} The randomised trials of Housing First find large gains in
housing stability and mixed to null effects on drug use, and they follow participants for
at most three years. Only one small trial, Milby et al.'s, of homeless adults already in
treatment, compares a sobriety condition with unconditional housing, which is the comparison
the paper is about, so it is the trial the model is calibrated to and checked against
(Sections 1 and 2).

\textbf{The price elasticities.} The model's participation elasticity is $-1.32$, somewhat
larger in absolute value than the $-0.3$ to $-0.9$ in Saffer and Chaloupka (1999) and
Grossman and Chaloupka (1998), and its intensive elasticity is $-0.63$, in the middle of the
$-0.3$ to $-0.8$ in Grossman and Chaloupka (1998) and Olmstead et al.\ (2015). A larger
price response works against Treatment First, since the screen lowers the drug price and a
cheaper drug draws users back, so in this respect the welfare gain is conservative (Section
5.2, Block 5).

In the paper: Section 4.6; Section 5.1, Blocks 2 and 4, and Table 3; Section 6, ``The
screen'', and footnotes 33 to 35; Section 5.2, Blocks 5, 7 and 8, and Table 4; Sections 1
and 2.
\end{response}

\cmt{C2}{A large homeless population is a negative externality}
\begin{referee}
``\dots a large homeless population, especially if many are addicted to drugs,
creates a negative externality beyond the fiscal burden of providing services.
\dots One structural rental housing paper that considers the negative externality
of homelessness (though doesn't have drug addiction) is Corbae, et al (2024).
\dots the externality seems very easy to introduce into the current model (though
would require a new calibration moment).''
\end{referee}
\begin{response}
\open I solved the model with this disamenity, but the welfare numbers in the paper
leave it out. I priced it from a direct valuation rather than from a new calibration
moment. In a contingent valuation, Loubi\`ere et al.\ (2020) find that Europeans would pay
about 0.2 percent of average earnings a year to end homelessness. With a disamenity of that
size, no long run gain in Table 5 moves by more than 0.14 percentage points. Because it
barely changes the results, I report it in a footnote and do not carry it in the benchmark.

On your point that the externality is larger when many of the homeless have a disorder:
under Treatment First the homeless population is both smaller and less often has a SUD. In
the main experiment, homelessness falls from 0.62 to 0.55 percent of households, and the
share of the homeless with a SUD falls from 20.2 to 15.6 percent (Table 5).

The paper does not cite Corbae, Glover and Nattinger (2024). The size of the disamenity in
the footnote comes from Loubi\`ere et al.'s valuation instead.

The welfare comparison turns instead on the children of users, which I take up at C3.

In the paper: Section 2, footnote 4; Section 6.2 and Table 5.
\end{response}

\cmt{C3}{Parents suffer anxiety about their children}
\begin{referee}
``\dots parents suffer anxiety that their children may become drug addicts and/or
homeless in adulthood. \dots the dynastic element seems like more of a stretch,
but could be discussed (I think it would only amplify the welfare gains from
reducing addiction).''
\end{referee}
\begin{response}
\newresult The link between parent and child that you describe is now in the model, and
the children of users are what the paper is about; that is why the title changed. You and
the Associate Editor, who objected that the previous version discussed intergenerational
transmission without any link between generations in the model, pointed to the same
missing piece. The model differs from your framing in one respect that matters for the
results. The parent does care about its child, but only in part, and most of the gain from
Treatment First is the part of the child's welfare that the parent does not count.

\textbf{How the parent values its child.} Every household has a child, who lives with it
for the parent's first 25 years, from 26 to 50, or until the parent dies if that comes
first. The child shares the household's consumption and dwelling but not the drug, and then
forms a household of its own, with a child of its own (Sections 4.1 and 4.2; Appendix A,
Block 9). The parent values its child at the weight $\omega_c$ in two places (equation
(24)): the child's flow utility while it lives at home, and the child's adult life, through
the loss the parent's own disorder inflicts on it. The second term is a difference: the
child's value at entry to adult life had the parent never had a SUD, less its value given
the parent's actual history (equation (25)). It is zero for a parent who never had a SUD and
positive otherwise. Written as a level, the term would be worth more to the households
likelier to die, and so most to those with a SUD, for reasons that have nothing to do with
the child (footnote 23).

\textbf{What the child inherits.} You thought the dynastic element would be a stretch. In
the model it takes the one weight and the two terms above, and a third value of the
disorder, so that the parent's charge can depend on whether it ever had a SUD (Section
4.1). Every entrant to adult life starts with the same endowment, in work, without a SUD,
without a voucher, and with SMI at the prevalence observed at age 26 (Section 4.13). What it
takes from its parent is its productivity, drawn around the parent's and lowered by 24
percent for life if the parent had a SUD at any point while the child lived at home
(Section 4.8, equation (22)). The 24 percent is the earnings loss Nilsson (2017) measures
for children exposed to alcohol in the womb in a Swedish experiment. The paper uses it as
the closest causal estimate, notes that this shock does not map one to one into a parent's
disorder after birth, and checks its size through the transmission of the SUD, below
(Section 2; Section 5.1, Block 5). The lower draw makes the child likelier to be poor, to
use, to acquire a SUD and SMI, to be homeless and to die earlier. The two outcomes you
name, a child who becomes addicted or homeless as an adult, therefore enter the parent's
problem through that term, at the weight $\omega_c$, to the extent that the parent's own
disorder changes them.

\textbf{How much the parent cares.} Altonji et al.\ (1997) reject perfect altruism, but
their test needs transfers between a parent and its grown child, which the model does not
have. So $\omega_c$ is set by behaviour. In the 2024 National Survey on Drug Use and
Health (NSDUH), an unmarried adult living with one other adult and no child was 17.5
percent more likely to have used an illicit drug in the past year than a lone parent living
with one child, with the lone parents reweighted to the other group's age, sex, marital
status and family income. Every household in the model has a child, so the model's
counterpart is a counterfactual: $\omega_c$ is the weight at which giving the child no
weight raises drug use by as much in proportion, 17.7 percent, and the calibrated weight is
0.085 (Section 5.1, Block 5, Table 3, and footnote 29). With the weight Nishiyama (2000)
estimates, 0.69, at the upper end of the literature, and with $b$ matched again to drug
use, Treatment First still raises a newborn's welfare, by 3.18 percent of lifetime
consumption against 2.57. Less of the gain is then an externality, 23 percent of it against
76 percent (footnote 28).

\textbf{Three externalities.} Because the parent counts only the share $\omega_c$ of what
its choices cost its child, the rest is an externality, and the paper names three (Section
1). Two are costs of a parent's drug use that Housing First leaves in place: the
consumption and housing the child goes without when the parent spends on the drug or a SUD
keeps it out of work (the budget externality), and the lower productivity the child starts
adult life with if the parent had a SUD while it lived at home (the earnings externality).
The third is the one many believe Treatment First creates: a child lives wherever its
parent lives, so a parent who loses its voucher and ends up in a shelter takes its child
there (the childhood homelessness externality). The welfare criterion is the newborn's, and
it counts the child's childhood at its full weight (Section 6.1).

\textbf{What it does to the answer.} You expected the link between generations to amplify
the welfare gains from reducing addiction. In this model it decides their sign. Without the
externalities, the child appears only as an adult: its parent gives it no weight, its
productivity does not depend on its parent's disorder, and its childhood is not counted. A
newborn then loses under Treatment First, by 1.17 percent of lifetime consumption with the
drug price and the chance of drawing a voucher held, 1.09 percent with that chance
balancing the budget, and 0.74 percent with both adjusting. With the externalities it gains
2.82, 6.45 and 2.57 percent in the same three cases (Section 6.2, Table 5). In the main
experiment the 2.57 percent splits into 1.50 percent from childhood and 0.46 percent from
better productivity at entry, both of which the parent does not count, 0.43 percent from the
newborn's own adult life under the policy, and 0.19 percent that the parent does count
(Section 6.3, Table 6). The parent's own concern for its child, the margin you raised, is
therefore the smallest of the four parts. The earnings externality is 18 percent of the
gain, and Treatment First still raises welfare with the penalty set to zero (footnotes 5
and 30).

\textbf{Childhood homelessness.} In the long run, the externality that critics of
Treatment First fear runs the other way. In the main experiment the share of children in a
shelter falls from 1.03 to 0.92 percent, and along the main transition, in which a
household with a SUD keeps its voucher by entering treatment and not using, it never rises
(Sections 6.2 and 6.4). It does rise in the first years when households with a SUD cannot
stop using: with no exemption and the calibrated detection rate, by up to 3.1 percent, and
it is above its Housing First level in periods 2 to 11 (Section 6.5, Table 7). The paper
also notes that the model likely overstates childhood homelessness, since every homeless
parent in the model takes its child into the shelter (footnote 7).

\textbf{Whether users have children.} They do. In the National Survey of Family Growth
2017 to 2019, 75 percent of adults aged 40 to 49 who used hard drugs in the past year had a
child, against 81 percent of those who used none, and among adults aged 26 to 49, 40
percent of those with a drug use disorder have a child under 18 at home, against 48 percent
of comparable adults who used no illicit drug. Even if all of the gain that comes in
childhood shrank in proportion to that gap, the gain would be only 11 percent smaller
(Section 2). In the model, 5.82 percent of children live with a parent who has a SUD,
against 4.6 percent in NSDUH 2024 (Section 5.2, Block 6).

\textbf{Transmission.} In the model, a child of a parent with a SUD has 1.46 times the odds
of acquiring one, at the low end of the 1.44 to 2.51 in Mellentin et al.\ (2016) and Yoon
et al.\ (2013), which suggests that the penalty is at the low end too. A homeless adult has
6.23 times the odds of being the child of someone who was homeless as well, above the 1.8
to 2.4 in Bassuk et al.\ (1997) (Section 5.2, Block 6).

In the paper: Sections 1 and 2; Sections 4.1, 4.2, 4.8, 4.10 and 4.13, and footnote 23;
Section 5.1, Block 5, Table 3, and footnotes 28 and 29; Section 5.2, Block 6; Sections 6.1
to 6.5, Tables 5 to 7, and footnotes 5, 7 and 30; Appendix A, Block 9.
\end{response}

\cmt{C4}{The rich model is a black box; show the policy functions}
\begin{referee}
``\dots it can be a bit of a black box. It would be helpful to have a section
where the mechanisms of the model are explored more thoroughly. For example, it
would be useful to show and explain intuitively a plot of the drug use policy
function for addicts and non-addicts as their income and wealth vary.''
\end{referee}
\begin{response}
\open The paper does not include the plot you describe, of the drug use policy function
by income and wealth for households with and without a SUD. In the quantitative model the
decision to use depends on age, assets, productivity, employment, the disorder, SMI and
the voucher, so I chose to show the mechanism where it can be isolated, in a stylized
model, and to show the quantitative model's drug use by income against the data. The
paper explores the mechanisms in three ways.

\textbf{A stylized model.} Section 3 is a new two period model, calibrated (Tables 1 and
2), whose sole purpose is to show why drug use concentrates at the bottom of the income
distribution. Households there differ only in productivity, so each consequence of a SUD
can be scaled up and down with the other switched off, and its effect on who uses read off
directly (Figure 1, with $b$ recalibrated at every point so that the share who use stays at
8.90 percent). With no consequences, the average user would earn 3.51 times the population
average. The risk of dying alone takes that to 0.64 at its calibrated strength. The risk of
losing one's job alone leaves it at 3.51, because saving partly insures it; with saving
switched off, the job risk alone takes it to 0.75. Section 3.3 gives the intuition, which I
set out at C5.

\textbf{Drug use by income in the quantitative model.} Figure 2 is the closest the paper
comes to your plot: the share of households that use drugs at each level of income
relative to the mean, model against NSDUH 2024, pooled over the other states. The
calibration targets only the average income of users, so the profile is untargeted. As in
the data, use is highest at the bottom and lowest at the top: 18.3 percent of households
below 0.22 of mean income use, against 13.7 percent in the data, and 4.2 percent of those
above 1.62 times the mean, against 4.8 percent. Between 0.43 and 1.62 times the mean, use
in the model is flatter and lower than in the data, 5.9 to 6.3 percent against 7.2 to 10.2.
Both differences lower users' average income in the model, which is why it falls slightly
short of its target, 62 percent of the average against 74 (Section 5.2, Block 2, and Table
3).

\textbf{The experiments, one step at a time.} Table 5 introduces Treatment First one
equilibrium condition at a time, with and without the externalities on children: the
screen alone, then with the chance of drawing a voucher balancing the budget, then with the
drug price clearing, and then with both. Section 6.3 and Table 6 split the newborn's gain
into four parts, two of them externalities, and say which the parent counts.

In the paper: Section 3 and Figure 1; Section 5.2, Block 2, and Figure 2; Sections 6.2
and 6.3, and Tables 5 and 6.
\end{response}

\cmt{C5}{The utility specification bakes in an income contingent gain or loss}
\begin{referee}
``It would also help to further explain the importance of the period utility
specification in (13) which, I think, implies that drug addiction is more
attractive for low income people \dots it also bakes in an income-contingent
gain/loss from becoming addicted by shifting the utility weight towards drug
consumption away from non-drug consumption and housing.''
\end{referee}
\begin{response}
\concede You and Referee~3, in comment M2(a), identified the same problem with the previous
specification. The weight on the drug rose when a household acquired the disorder, so
reducing the disorder mechanically raised the utility of consumption and housing, and a
policy that reduced it gained welfare through the utility function itself, not only through
behaviour.

In the revised model, preferences are the same for every household. There is one weight on
the drug, $\theta$, one weight on housing, $\gamma$, and one shifter, $b$, whatever a
household's SUD status, its mental health or its dwelling (Section 4.2, equations (11) and
(12)), and SMI does not enter utility at all (Section 4.7). The Associate Editor suggested
keeping the weights fixed and letting $b$ vary with the disorder; the revision keeps $b$
fixed as well. In period utility, the only term tied to the disorder is the cost of a year
in treatment, which a household with a SUD chooses whether to pay. What a SUD costs runs
through the risks it brings and through the child: a higher chance of dying, a lower chance
of finding work, a higher chance of SMI, less consumption and housing for the child when the
parent spends on the drug or is out of work, and a productivity penalty on the child, who
inherits a worse draw, not the disorder or a taste for the drug (Sections 4.6 to 4.8).

On your question, whether the utility function makes drug use more attractive at low
income: in the revised model it does the opposite. Drugs are a luxury good bought at a
price. Because $b>0$, the first unit has a finite marginal value, so most households buy
none, and a household with more income gives up less utility for each dollar it spends on
the drug. With no consequences of use, use would therefore concentrate at the top: in the
stylized model, the average user would earn 3.51 times the population average (Section
3.3).

What reverses the gradient is the risk of dying, which cannot be insured. Dying takes away
the whole continuation value of being alive, which grows with income, so it deters high
income households and not low income ones. Losing a job is partly insured by saving, so
its cost is roughly proportional to income and, on its own, does not change who uses. In
the stylized model the risk of dying alone takes the average user from 3.51 times the
population average to 0.64, and with both risks at their calibrated strength it is 0.74,
the NSDUH 2024 value the stylized model is calibrated to (Figure 1 and footnote 14). So the
cost of a SUD does depend on income, as you suspected, but through the mortality risk of
the disorder rather than through preferences, and it is larger for richer households. In
the quantitative model the mortality multiplier of a SUD, 4.00, is calibrated to the income
of users and lies within the 3.91 to 6.6 range of mortality rate ratios in the literature;
users' income falls slightly short of its target, 62 percent of the average against 74
(Section 5.1, Block 1). Figure 2 then compares the whole profile of use by income, which is
not targeted, with the data (C4).

Keeping preferences fixed is also what lets the paper explain why users are poor without
assuming that preferences differ systematically across households. Greenwood et al.\
(2025) reproduce the same fact by letting preferences for the drug differ by education;
here no preference parameter varies with income (Section 2).

In the paper: Sections 4.2 and 4.7; Section 3.3, Figure 1 and footnote 14; Section 5.1,
Block 1; Section 5.2, Block 2, and Figure 2; Section 2.
\end{response}

\cmt{C6}{Optimal voucher design, and $\chi_H$ looks high}
\begin{referee}
``The parameter $\chi_H$ governs the fiscal cost of having somebody unhoused and
appears pretty high at 25\% of average income. Would there be significant further
welfare gains from broadening the voucher program by making higher incomes
eligible or covering more of the rent (with the drug-free contingency)? My
(perhaps wrong) reading is that the cost of having somebody unhoused is
significantly higher than the rent on the smallest housing unit.''
\end{referee}
\begin{response}
\open Two things have changed, and the second means the paper does not answer your
question directly.

\textbf{The fiscal cost of homelessness.} The revised model has no fiscal cost of
homelessness. The shelter is free, and the government's budget carries the income taxes,
the estates of the dead, the endowment of entrants, unemployment insurance, social
security, food assistance, the voucher, the drug test and treatment, and nothing for a
household in a shelter (Section 4.11, equation (26)). So the parameter you asked about is
gone, and with it the comparison between the cost of a homeless household and the rent on
the smallest unit. The paper does not solve the model with such a cost. Since Treatment
First lowers homelessness in the long run and does not raise it along the main transition,
I would expect a cost per homeless household to add to what Treatment First saves rather
than to what it spends. What pays for more vouchers under Treatment First is instead the
vouchers revoked from users who are caught, and higher tax revenue, mainly because more
adults start adult life without the earnings penalty and to a lesser extent because fewer
are out of work with a SUD. In the main experiment the probability of drawing a voucher
rises from 6.75 to 7.80 percent (Section 6.2).

\textbf{Broader programmes.} The previous version compared programmes with and without
rationing. The revision does not, and it does not solve for the optimal programme. It keeps
the programme rationed and its rules at those of the Department of Housing and Urban
Development (HUD): a household is eligible below half of median income, and the voucher
pays the rent above 30 percent of the holder's income, up to the fair market rent, which is
the average rent paid by housed households in the benchmark (Section 4.3, equation (15);
Appendix A, Block 12). In the model a voucher holder pays 39.0 percent of its income in
rent, against 31.2 percent in HUD's data, because a holder that rents above the fair market
rent pays the difference itself (Section 5.2, Block 3, and footnote 32). Every experiment
holds the tax system at its benchmark level, and either holds net government revenue at its
benchmark level through the chance of drawing a voucher or leaves the savings unspent. It
asks only whether the assistance that exists should require sobriety (Sections 4.11 and 6).
Raising the eligibility cutoff or the share of the rent covered is not among the
experiments, so the paper cannot tell you whether either would add further gains.

What the paper does show is the value of more vouchers that carry the sobriety condition,
at the same net revenue. When the probability of drawing a voucher is free to spend what
the screen saves, with the drug price held, it rises from 6.75 to 9.24 percent,
homelessness falls by 20 percent and the share of children in a shelter by 19 percent, and
a newborn gains 6.45 percent of lifetime consumption, against 2.82 percent when the savings
go unspent. The parents gain 2.42 percent of their own lifetime consumption, against a loss
of 1.19 percent (Section 6.2, Table 5, seventh and eighth columns; Table 6).

In the paper: Sections 4.3 and 4.11; Section 5.2, Block 3, and footnote 32; Section 6;
Section 6.2, and Tables 5 and 6; Appendix A, Block 12.
\end{response}

%======================================================================
\section*{Detail concerns}

\cmt{D1}{The housing grid}
\begin{referee}
``\dots the housing grid. It has seven values but only three are calibrated \dots
The full grid is then listed in footnote 4 rather than a calibration table. My
assumption before seeing the grid was that the bounds were calibrated and then the
sizes were distributed uniformly \dots But they are not uniformly spread out, at
least in levels or ratios. So, how are the housing levels being chosen and does it
matter for results? Could they just be spread uniformly without affecting the
calibration and policy conclusions? If not then I think they need to be included as
parameters in the calibration.''
\end{referee}
\begin{response}
\concede \open You were right that the version you read gave no reason for the spacing of
its grid. The revision makes the grid smaller and puts every point of it in a calibration
table, but it still gives no source or moment for the two upper rungs, and it does not test
whether their spacing matters.

\textbf{The grid is now four points.} Homeownership is gone, along with house prices,
endogenous rents and labour supply, which in earlier versions had a negligible effect on
the results (footnote 19). A household chooses between a shelter, $h_1$, and three rental
sizes, $h_2<h_3<h_4$, at a fixed rent per unit (Sections 4.4 and 4.5).

\textbf{The two points at the bottom are calibrated.} They decide what a voucher, and its
removal, are worth to a poor household. The shelter's service flow $h_1$ is set so that
homelessness is within 0.40 to 1.00 percent of households, the floor from HUD's 2023 point
in time count of a single night and the ceiling from the share homeless at some point over
a year in Burt et al.\ (2001). The model gives 0.62 percent, with $h_1$ at 53 percent of the
smallest rental (Section 5.1, Block 2, and Table 3). The smallest rental $h_2$ decides how
much of its income a poor household must spend to stay out of the shelter. It is set so
that the share of renters paying more than half of their income for housing is close to
the data: 19.2 percent in the 2022 Survey of Consumer Finances, against 25.8 percent in the
model (Appendix A, Block 8).

\textbf{The two upper rungs are set.} They are at 4.00 and 8.00 times the smallest rental,
and they now appear as parameters in Table A1 rather than in a footnote, marked as set,
with no moment or source. They are not uniform in levels or in ratios. They decide what a
household that can afford more than the smallest unit can rent. The paper places them among
the standard parameters, which it says matter less for the results than those of Table 3
(Appendix A), but it does not show this for the rungs, and it reports no results with a
different spacing. So the paper does not answer your last two questions, whether the
spacing matters and whether a uniform ladder would change the calibration or the
conclusions.

In the paper: Sections 4.4 and 4.5, and footnote 19; Section 5.1, Block 2, and Table 3;
Appendix A, Block 8, and Table A1.
\end{response}

\end{document}
